\section{Simulation design and complete results}

\subsection{Data-generating distributions}

All simulations were generated independently from ordinary finite mixtures,
but performance was evaluated about the corresponding population
classification maximum-likelihood target.  This distinction matters under
overlap.  For the two-component designs, the first component was multivariate
Normal with AR(1) covariance parameter 0.25.  Its location was
$-\delta\mathbf{e}/2$, where
$\mathbf{e}=p^{-1/2}(1,\ldots,1)^\mathsf T$.  The second location was
$\delta\mathbf{e}/2$.  We used $\delta=6.0,4.25,$ and $3.0$ for separated,
moderate, and strong overlap, respectively.  In dimensions $p\geq5$, the
strong-overlap distance was 3.75 because the prespecified $3.0$ design caused
population component collapse and therefore did not satisfy the regular-target
conditions of the inferential study.

In Normal--Student-$t$ experiments, the second component had 6 degrees of
freedom and an AR(1) scatter matrix with correlation 0.20 and marginal scale
1.1.  In Normal--skew-normal experiments, its covariance had AR(1) correlation
$-0.15$ and its shape vector was
$(3,-2,0,\ldots,0)^\mathsf T$.  Balanced weights were $(0.5,0.5)$ and the
imbalanced experiments used $(0.2,0.8)$.  The three-component design combined
Normal, Student-$t_7$, and skew-normal components with weights
$(0.3,0.4,0.3)$.  Their first two location coordinates formed an equilateral
triangle of radius 2.3; remaining location coordinates were zero.  Its shape
vector was $(2.5,-1.5,0,\ldots,0)^\mathsf T$.

The primary design comprised dimensions $p=2$ and $p=5$, sample sizes 500 and
2000, the three overlap regimes, imbalance at moderate overlap, two
three-component experiments at $n=1000$, and dimension-$10$ stress experiments
at $n=1000$.  There were 1000 replications in each principal cell and 500 in
each dimension-$10$ cell, giving 31,000 attempted primary fits.  An additional
large-sample extension used 500 replications at $n=8000$ for the two
two-dimensional strong-overlap models and at $n=4000$ for the two
three-component models.

\subsection{Targets, estimators, and performance measures}

For each of the 18 population scenarios, we generated ten independent
reference samples of size 100,000; five were used in the initial dimension-10
stress computation.  All reference fits succeeded.  The selected population
target was the coordinate mean of the canonicalized reference solutions.
Across scenarios, the largest coordinate Monte Carlo standard error was 0.033
and the largest range across reference estimates was 0.322.  Target uncertainty
was retained in the archive and audited before the confirmatory run.

Each sample was fitted using the multistart and family-assignment procedure in
the main paper.  For every successful fit, inference was evaluated with
bandwidth multipliers $0.75$, 1, and 1.5 around the pilot bandwidth.  The
primary corrected method is the stabilized estimator at multiplier 1.  A
coordinate contributed to coverage only if the base estimator and the
corresponding covariance calculation were available.  We report the number of
usable replications and do not replace failed seeds.  Scenario summaries are
medians across unconstrained parameter coordinates; mean coordinate coverage
and all coordinate-level results remain in the machine-readable archive.

Table~\ref{tab:s-full-performance} gives every scenario-level primary and
large-sample result for naive and stabilized inference.  It shows that the
correction is most important for Normal--$t$ overlap and imbalance, and that
several difficult skew-normal cells converge slowly at $n=500$.  The
large-sample extension resolves much of this discrepancy, except for the
two-dimensional three-component experiment, where family and face
interactions remain difficult.  This cell was retained as a limitation.

Figure~\ref{fig:s-bandwidth} gives the full bandwidth sensitivity display.
Thin lines are individual scenario--sample-size cells and the black line is
their median. Aggregate conclusions are stable across the prespecified
multipliers, although individual difficult cells remain sensitive.

\begin{figure}[t]
\centering
\includegraphics[width=\textwidth]{figures/bandwidth_sensitivity.pdf}
\caption{Sensitivity of stabilized inference to bandwidth multipliers 0.75,
1, and 1.5. Thin coloured lines are individual cells; black is the across-cell
median.}
\label{fig:s-bandwidth}
\end{figure}

\subsection{Failure and stabilization accounting}

Table~\ref{tab:s-failures} reports all primary fit failures and all
replications in which at least one raw bandwidth-specific covariance was
nonpositive.  Normal--Student-$t$ fitting succeeded in every primary
replication.  Skew-normal fit failure ranged from 0 to 7.6\%, with the largest
rate in the dimension-10 cell.  Raw curvature failure is appreciably more
common and increases under overlap and dimension.  These are failures of the
finite-sample plug-in curvature, not evidence that the population curvature is
nonpositive.  The stabilization makes the covariance computable while its
activation frequency generally decreases with sample size, as predicted by
the asymptotic-inactivity proposition.

\subsection{Direct convergence experiment for the boundary estimator}

We separately estimated a scalar functional of the boundary curvature at
sample sizes 250, 1000, 4000, and 16,000.  The results were
\begin{center}
\begin{tabular}{rrrrr}
\toprule
$n$ & Bias & Empirical SD & RMSE & Bandwidth\\
\midrule
250   & $-0.00223$ & 0.04519 & 0.04520 & 0.33145\\
1000  & $-0.00432$ & 0.02903 & 0.02932 & 0.25119\\
4000  & $-0.00236$ & 0.01678 & 0.01693 & 0.19037\\
16000 & $-0.00077$ & 0.00958 & 0.00960 & 0.14427\\
\bottomrule
\end{tabular}
\end{center}
The shrinking bias and variance provide a direct numerical check of the
kernel-estimator consistency result, separate from Wald-interval performance.

\subsection{Direct evidence for the limiting theory}

Figure~\ref{fig:s-asymptotic-evidence} evaluates the consistency, root-$n$
rate, and Gaussian-limit claims using the replication-level estimates, rather
than treating interval coverage as a proxy. For every successful replication
we formed $\sqrt n(\widehat{\boldsymbol\theta}-\boldsymbol\theta^\star)$ in
the identifiable unconstrained coordinates. The upper-left panel compares
coordinate-standardized empirical quantiles with standard-normal quantiles;
the upper-right reports the coordinate-level Kolmogorov distance after
centering and scaling. The lower panels show $\sqrt n$ times coordinate-median
RMSE and absolute bias for every scenario observed at both $n=500$ and
$n=2000$. These are diagnostics, not formal proofs: their purpose is to show
whether the finite-sample experiments display the behaviour asserted by the
theorems and to retain deviations in difficult cells.

\begin{figure}[t]
\centering
\includegraphics[width=\textwidth]{figures/asymptotic_evidence.pdf}
\caption{Empirical evidence for the limiting theory. Top left: median
coordinate-level quantiles of standardized root-$n$ errors. Top right:
coordinate-level Kolmogorov distance from a fitted standard normal law.
Bottom: scenario trajectories of scaled RMSE and scaled absolute bias.}
\label{fig:s-asymptotic-evidence}
\end{figure}

\subsection{Complete coordinate-level diagnostic atlas}

Figures~\ref{fig:s-bias-atlas}--\ref{fig:s-se-atlas} retain every
unconstrained coordinate from every confirmatory scenario. They are the direct
counterparts of the parameter-by-parameter bias, RMSE, coverage, and
SD-versus-SE plots in the previous study, but use the redesigned targets,
families, dimensions, sample sizes, and stabilization. A separate panel is
shown for each scenario. Colour denotes sample size; in the inference atlases,
dotted and solid curves denote naive and stabilized boundary-aware inference.
The machine-readable source values accompany the figures.

\begin{figure}[p]\centering
\includegraphics[width=.94\textwidth]{figures/coordinate_bias_atlas.pdf}
\caption{Coordinate-level bias atlas for all confirmatory scenarios.}
\label{fig:s-bias-atlas}\end{figure}
\begin{figure}[p]\centering
\includegraphics[width=.94\textwidth]{figures/coordinate_rmse_atlas.pdf}
\caption{Coordinate-level RMSE atlas for all confirmatory scenarios.}
\label{fig:s-rmse-atlas}\end{figure}
\begin{figure}[p]\centering
\includegraphics[width=.94\textwidth]{figures/coordinate_coverage_atlas.pdf}
\caption{Coordinate-level coverage atlas for naive and stabilized inference.}
\label{fig:s-coverage-atlas}\end{figure}
\begin{figure}[p]\centering
\includegraphics[width=.94\textwidth]{figures/coordinate_se_atlas.pdf}
\caption{Coordinate-level estimated-SE to empirical-SD ratio atlas.}
\label{fig:s-se-atlas}\end{figure}

\section{Application screening and additional results}

The prespecified screen considered breast cancer, wine, digits 0 versus 1,
and three Dry Bean pairs; Normal--Student-$t$ and Normal--skew-normal fits were
examined after two- and five-dimensional principal-component reduction.
Table~\ref{tab:s-screen} reports the complete screen.  Fits with negligible
kernel support were not treated as useful evidence for a boundary correction,
and failed covariance calculations were retained.

For the selected Dry Bean analyses, Tables~\ref{tab:s-dermason-parameters} and
\ref{tab:s-cali-parameters} report every unconstrained coordinate.  The local
bootstrap uses 500 resamples and the unrestricted diagnostic uses 200.  In
Dermason--Seker, stabilized standard errors track local-bootstrap variability
across the full parameter vector.  Cali--Barbunya shows a consistent movement
in the correct direction but incomplete calibration and a high local
optimization failure rate.  The unrestricted refits have very large
location-coordinate variability because they can select another basin; these
global standard deviations are retained in the CSV archive but are not pooled
with the local Gaussian target.

\begingroup\scriptsize\setlength{\tabcolsep}{2pt}
\begin{longtable}{p{0.27\linewidth}rrlrrr}
\caption{Complete scenario-level simulation summary. SE/SD and coverage are coordinate medians; stabilization is the fraction of usable fits for which the eigenvalue constraint was active.}\label{tab:s-full-performance}\\
\toprule
Scenario & $n$ & Method & Usable & SE/SD & Coverage & Stabilization\\
\midrule\endfirsthead
\toprule
Scenario & $n$ & Method & Usable & SE/SD & Coverage & Stabilization\\
\midrule\endhead
\multicolumn{7}{l}{\textit{Primary study}}\\
nsn\_\allowbreak{}p10\_\allowbreak{}moderate\_\allowbreak{}balanced & 1000 & Corrected & 462 & 1.030 & 0.942 & 0.346 \\
nsn\_\allowbreak{}p10\_\allowbreak{}moderate\_\allowbreak{}balanced & 1000 & Naive & 462 & 0.970 & 0.937 & 0.000 \\
nsn\_\allowbreak{}p2\_\allowbreak{}moderate\_\allowbreak{}balanced & 500 & Corrected & 993 & 0.565 & 0.872 & 0.085 \\
nsn\_\allowbreak{}p2\_\allowbreak{}moderate\_\allowbreak{}balanced & 500 & Naive & 993 & 0.517 & 0.863 & 0.000 \\
nsn\_\allowbreak{}p2\_\allowbreak{}moderate\_\allowbreak{}balanced & 2000 & Corrected & 998 & 0.748 & 0.937 & 0.029 \\
nsn\_\allowbreak{}p2\_\allowbreak{}moderate\_\allowbreak{}balanced & 2000 & Naive & 998 & 0.650 & 0.923 & 0.000 \\
nsn\_\allowbreak{}p2\_\allowbreak{}moderate\_\allowbreak{}imbalanced & 500 & Corrected & 993 & 0.527 & 0.889 & 0.063 \\
nsn\_\allowbreak{}p2\_\allowbreak{}moderate\_\allowbreak{}imbalanced & 500 & Naive & 993 & 0.502 & 0.880 & 0.000 \\
nsn\_\allowbreak{}p2\_\allowbreak{}moderate\_\allowbreak{}imbalanced & 2000 & Corrected & 1000 & 0.805 & 0.946 & 0.027 \\
nsn\_\allowbreak{}p2\_\allowbreak{}moderate\_\allowbreak{}imbalanced & 2000 & Naive & 1000 & 0.738 & 0.939 & 0.000 \\
nsn\_\allowbreak{}p2\_\allowbreak{}separated\_\allowbreak{}balanced & 500 & Corrected & 984 & 0.557 & 0.889 & 0.024 \\
nsn\_\allowbreak{}p2\_\allowbreak{}separated\_\allowbreak{}balanced & 500 & Naive & 984 & 0.535 & 0.888 & 0.000 \\
nsn\_\allowbreak{}p2\_\allowbreak{}separated\_\allowbreak{}balanced & 2000 & Corrected & 997 & 0.975 & 0.945 & 0.006 \\
nsn\_\allowbreak{}p2\_\allowbreak{}separated\_\allowbreak{}balanced & 2000 & Naive & 997 & 0.966 & 0.943 & 0.000 \\
nsn\_\allowbreak{}p2\_\allowbreak{}strong\_\allowbreak{}balanced & 500 & Corrected & 997 & 0.613 & 0.776 & 0.250 \\
nsn\_\allowbreak{}p2\_\allowbreak{}strong\_\allowbreak{}balanced & 500 & Naive & 997 & 0.484 & 0.752 & 0.000 \\
nsn\_\allowbreak{}p2\_\allowbreak{}strong\_\allowbreak{}balanced & 2000 & Corrected & 1000 & 0.762 & 0.922 & 0.138 \\
nsn\_\allowbreak{}p2\_\allowbreak{}strong\_\allowbreak{}balanced & 2000 & Naive & 1000 & 0.573 & 0.868 & 0.000 \\
nsn\_\allowbreak{}p5\_\allowbreak{}moderate\_\allowbreak{}balanced & 500 & Corrected & 936 & 0.978 & 0.904 & 0.197 \\
nsn\_\allowbreak{}p5\_\allowbreak{}moderate\_\allowbreak{}balanced & 500 & Naive & 936 & 0.934 & 0.894 & 0.000 \\
nsn\_\allowbreak{}p5\_\allowbreak{}moderate\_\allowbreak{}balanced & 2000 & Corrected & 989 & 0.994 & 0.943 & 0.119 \\
nsn\_\allowbreak{}p5\_\allowbreak{}moderate\_\allowbreak{}balanced & 2000 & Naive & 989 & 0.892 & 0.933 & 0.000 \\
nsn\_\allowbreak{}p5\_\allowbreak{}separated\_\allowbreak{}balanced & 500 & Corrected & 941 & 0.927 & 0.905 & 0.039 \\
nsn\_\allowbreak{}p5\_\allowbreak{}separated\_\allowbreak{}balanced & 500 & Naive & 941 & 0.914 & 0.904 & 0.000 \\
nsn\_\allowbreak{}p5\_\allowbreak{}separated\_\allowbreak{}balanced & 2000 & Corrected & 987 & 0.977 & 0.946 & 0.018 \\
nsn\_\allowbreak{}p5\_\allowbreak{}separated\_\allowbreak{}balanced & 2000 & Naive & 987 & 0.969 & 0.945 & 0.000 \\
nsn\_\allowbreak{}p5\_\allowbreak{}strong\_\allowbreak{}balanced & 500 & Corrected & 954 & 1.515 & 0.885 & 0.283 \\
nsn\_\allowbreak{}p5\_\allowbreak{}strong\_\allowbreak{}balanced & 500 & Naive & 954 & 0.905 & 0.868 & 0.000 \\
nsn\_\allowbreak{}p5\_\allowbreak{}strong\_\allowbreak{}balanced & 2000 & Corrected & 987 & 1.003 & 0.944 & 0.173 \\
nsn\_\allowbreak{}p5\_\allowbreak{}strong\_\allowbreak{}balanced & 2000 & Naive & 987 & 0.872 & 0.928 & 0.000 \\
nt\_\allowbreak{}p10\_\allowbreak{}moderate\_\allowbreak{}balanced & 1000 & Corrected & 500 & 1.000 & 0.942 & 0.800 \\
nt\_\allowbreak{}p10\_\allowbreak{}moderate\_\allowbreak{}balanced & 1000 & Naive & 500 & 0.878 & 0.920 & 0.000 \\
nt\_\allowbreak{}p2\_\allowbreak{}moderate\_\allowbreak{}balanced & 500 & Corrected & 1000 & 0.920 & 0.924 & 0.105 \\
nt\_\allowbreak{}p2\_\allowbreak{}moderate\_\allowbreak{}balanced & 500 & Naive & 1000 & 0.758 & 0.885 & 0.000 \\
nt\_\allowbreak{}p2\_\allowbreak{}moderate\_\allowbreak{}balanced & 2000 & Corrected & 1000 & 0.994 & 0.940 & 0.034 \\
nt\_\allowbreak{}p2\_\allowbreak{}moderate\_\allowbreak{}balanced & 2000 & Naive & 1000 & 0.827 & 0.895 & 0.000 \\
nt\_\allowbreak{}p2\_\allowbreak{}moderate\_\allowbreak{}imbalanced & 500 & Corrected & 1000 & 0.844 & 0.895 & 0.217 \\
nt\_\allowbreak{}p2\_\allowbreak{}moderate\_\allowbreak{}imbalanced & 500 & Naive & 1000 & 0.615 & 0.834 & 0.000 \\
nt\_\allowbreak{}p2\_\allowbreak{}moderate\_\allowbreak{}imbalanced & 2000 & Corrected & 1000 & 1.009 & 0.931 & 0.104 \\
nt\_\allowbreak{}p2\_\allowbreak{}moderate\_\allowbreak{}imbalanced & 2000 & Naive & 1000 & 0.760 & 0.865 & 0.000 \\
nt\_\allowbreak{}p2\_\allowbreak{}separated\_\allowbreak{}balanced & 500 & Corrected & 1000 & 0.977 & 0.941 & 0.051 \\
nt\_\allowbreak{}p2\_\allowbreak{}separated\_\allowbreak{}balanced & 500 & Naive & 1000 & 0.884 & 0.927 & 0.000 \\
nt\_\allowbreak{}p2\_\allowbreak{}separated\_\allowbreak{}balanced & 2000 & Corrected & 1000 & 1.004 & 0.948 & 0.010 \\
nt\_\allowbreak{}p2\_\allowbreak{}separated\_\allowbreak{}balanced & 2000 & Naive & 1000 & 0.945 & 0.938 & 0.000 \\
nt\_\allowbreak{}p2\_\allowbreak{}strong\_\allowbreak{}balanced & 500 & Corrected & 1000 & 0.675 & 0.843 & 0.328 \\
nt\_\allowbreak{}p2\_\allowbreak{}strong\_\allowbreak{}balanced & 500 & Naive & 1000 & 0.491 & 0.750 & 0.000 \\
nt\_\allowbreak{}p2\_\allowbreak{}strong\_\allowbreak{}balanced & 2000 & Corrected & 1000 & 0.910 & 0.891 & 0.272 \\
nt\_\allowbreak{}p2\_\allowbreak{}strong\_\allowbreak{}balanced & 2000 & Naive & 1000 & 0.554 & 0.742 & 0.000 \\
nt\_\allowbreak{}p5\_\allowbreak{}moderate\_\allowbreak{}balanced & 500 & Corrected & 1000 & 0.977 & 0.934 & 0.417 \\
nt\_\allowbreak{}p5\_\allowbreak{}moderate\_\allowbreak{}balanced & 500 & Naive & 1000 & 0.869 & 0.908 & 0.000 \\
nt\_\allowbreak{}p5\_\allowbreak{}moderate\_\allowbreak{}balanced & 2000 & Corrected & 1000 & 1.021 & 0.943 & 0.149 \\
nt\_\allowbreak{}p5\_\allowbreak{}moderate\_\allowbreak{}balanced & 2000 & Naive & 1000 & 0.847 & 0.904 & 0.000 \\
nt\_\allowbreak{}p5\_\allowbreak{}separated\_\allowbreak{}balanced & 500 & Corrected & 1000 & 1.014 & 0.945 & 0.169 \\
nt\_\allowbreak{}p5\_\allowbreak{}separated\_\allowbreak{}balanced & 500 & Naive & 1000 & 0.958 & 0.938 & 0.000 \\
nt\_\allowbreak{}p5\_\allowbreak{}separated\_\allowbreak{}balanced & 2000 & Corrected & 1000 & 1.018 & 0.950 & 0.045 \\
nt\_\allowbreak{}p5\_\allowbreak{}separated\_\allowbreak{}balanced & 2000 & Naive & 1000 & 0.956 & 0.939 & 0.000 \\
nt\_\allowbreak{}p5\_\allowbreak{}strong\_\allowbreak{}balanced & 500 & Corrected & 1000 & 0.622 & 0.913 & 0.492 \\
nt\_\allowbreak{}p5\_\allowbreak{}strong\_\allowbreak{}balanced & 500 & Naive & 1000 & 0.530 & 0.873 & 0.000 \\
nt\_\allowbreak{}p5\_\allowbreak{}strong\_\allowbreak{}balanced & 2000 & Corrected & 1000 & 1.036 & 0.940 & 0.295 \\
nt\_\allowbreak{}p5\_\allowbreak{}strong\_\allowbreak{}balanced & 2000 & Naive & 1000 & 0.800 & 0.886 & 0.000 \\
ntsn\_\allowbreak{}p2\_\allowbreak{}moderate\_\allowbreak{}balanced & 1000 & Corrected & 992 & 0.567 & 0.747 & 0.196 \\
ntsn\_\allowbreak{}p2\_\allowbreak{}moderate\_\allowbreak{}balanced & 1000 & Naive & 992 & 0.460 & 0.696 & 0.000 \\
ntsn\_\allowbreak{}p5\_\allowbreak{}moderate\_\allowbreak{}balanced & 1000 & Corrected & 956 & 1.010 & 0.921 & 0.542 \\
ntsn\_\allowbreak{}p5\_\allowbreak{}moderate\_\allowbreak{}balanced & 1000 & Naive & 956 & 0.879 & 0.901 & 0.000 \\
\multicolumn{7}{l}{\textit{Large-sample extension}}\\
nsn\_\allowbreak{}p2\_\allowbreak{}strong\_\allowbreak{}balanced & 8000 & Corrected & 500 & 0.857 & 0.940 & 0.030 \\
nsn\_\allowbreak{}p2\_\allowbreak{}strong\_\allowbreak{}balanced & 8000 & Naive & 500 & 0.627 & 0.862 & 0.000 \\
nt\_\allowbreak{}p2\_\allowbreak{}strong\_\allowbreak{}balanced & 8000 & Corrected & 500 & 1.049 & 0.923 & 0.220 \\
nt\_\allowbreak{}p2\_\allowbreak{}strong\_\allowbreak{}balanced & 8000 & Naive & 500 & 0.560 & 0.721 & 0.000 \\
ntsn\_\allowbreak{}p2\_\allowbreak{}moderate\_\allowbreak{}balanced & 4000 & Corrected & 490 & 0.666 & 0.863 & 0.053 \\
ntsn\_\allowbreak{}p2\_\allowbreak{}moderate\_\allowbreak{}balanced & 4000 & Naive & 490 & 0.400 & 0.816 & 0.000 \\
ntsn\_\allowbreak{}p5\_\allowbreak{}moderate\_\allowbreak{}balanced & 4000 & Corrected & 464 & 1.023 & 0.942 & 0.254 \\
ntsn\_\allowbreak{}p5\_\allowbreak{}moderate\_\allowbreak{}balanced & 4000 & Naive & 464 & 0.864 & 0.905 & 0.000 \\
\bottomrule
\end{longtable}
\endgroup

\begingroup\scriptsize\setlength{\tabcolsep}{3pt}
\begin{longtable}{p{0.43\linewidth}rrrr}
\caption{Complete primary-study failure accounting. Raw curvature failures count replications in which at least one bandwidth produced a nonpositive unmodified curvature.}\label{tab:s-failures}\\
\toprule
Scenario & $n$ & Planned & Fit failures & Raw failures\\
\midrule\endfirsthead
\toprule
Scenario & $n$ & Planned & Fit failures & Raw failures\\
\midrule\endhead
nsn\_\allowbreak{}p10\_\allowbreak{}moderate\_\allowbreak{}balanced & 1000 & 500 & 38 & 161 \\
nsn\_\allowbreak{}p2\_\allowbreak{}moderate\_\allowbreak{}balanced & 500 & 1000 & 7 & 78 \\
nsn\_\allowbreak{}p2\_\allowbreak{}moderate\_\allowbreak{}balanced & 2000 & 1000 & 2 & 29 \\
nsn\_\allowbreak{}p2\_\allowbreak{}moderate\_\allowbreak{}imbalanced & 500 & 1000 & 7 & 54 \\
nsn\_\allowbreak{}p2\_\allowbreak{}moderate\_\allowbreak{}imbalanced & 2000 & 1000 & 0 & 19 \\
nsn\_\allowbreak{}p2\_\allowbreak{}separated\_\allowbreak{}balanced & 500 & 1000 & 16 & 22 \\
nsn\_\allowbreak{}p2\_\allowbreak{}separated\_\allowbreak{}balanced & 2000 & 1000 & 3 & 8 \\
nsn\_\allowbreak{}p2\_\allowbreak{}strong\_\allowbreak{}balanced & 500 & 1000 & 3 & 179 \\
nsn\_\allowbreak{}p2\_\allowbreak{}strong\_\allowbreak{}balanced & 2000 & 1000 & 0 & 81 \\
nsn\_\allowbreak{}p5\_\allowbreak{}moderate\_\allowbreak{}balanced & 500 & 1000 & 64 & 180 \\
nsn\_\allowbreak{}p5\_\allowbreak{}moderate\_\allowbreak{}balanced & 2000 & 1000 & 11 & 109 \\
nsn\_\allowbreak{}p5\_\allowbreak{}separated\_\allowbreak{}balanced & 500 & 1000 & 59 & 32 \\
nsn\_\allowbreak{}p5\_\allowbreak{}separated\_\allowbreak{}balanced & 2000 & 1000 & 13 & 19 \\
nsn\_\allowbreak{}p5\_\allowbreak{}strong\_\allowbreak{}balanced & 500 & 1000 & 46 & 250 \\
nsn\_\allowbreak{}p5\_\allowbreak{}strong\_\allowbreak{}balanced & 2000 & 1000 & 13 & 164 \\
nt\_\allowbreak{}p10\_\allowbreak{}moderate\_\allowbreak{}balanced & 1000 & 500 & 0 & 372 \\
nt\_\allowbreak{}p2\_\allowbreak{}moderate\_\allowbreak{}balanced & 500 & 1000 & 0 & 73 \\
nt\_\allowbreak{}p2\_\allowbreak{}moderate\_\allowbreak{}balanced & 2000 & 1000 & 0 & 13 \\
nt\_\allowbreak{}p2\_\allowbreak{}moderate\_\allowbreak{}imbalanced & 500 & 1000 & 0 & 159 \\
nt\_\allowbreak{}p2\_\allowbreak{}moderate\_\allowbreak{}imbalanced & 2000 & 1000 & 0 & 79 \\
nt\_\allowbreak{}p2\_\allowbreak{}separated\_\allowbreak{}balanced & 500 & 1000 & 0 & 44 \\
nt\_\allowbreak{}p2\_\allowbreak{}separated\_\allowbreak{}balanced & 2000 & 1000 & 0 & 10 \\
nt\_\allowbreak{}p2\_\allowbreak{}strong\_\allowbreak{}balanced & 500 & 1000 & 0 & 142 \\
nt\_\allowbreak{}p2\_\allowbreak{}strong\_\allowbreak{}balanced & 2000 & 1000 & 0 & 81 \\
nt\_\allowbreak{}p5\_\allowbreak{}moderate\_\allowbreak{}balanced & 500 & 1000 & 0 & 346 \\
nt\_\allowbreak{}p5\_\allowbreak{}moderate\_\allowbreak{}balanced & 2000 & 1000 & 0 & 96 \\
nt\_\allowbreak{}p5\_\allowbreak{}separated\_\allowbreak{}balanced & 500 & 1000 & 0 & 154 \\
nt\_\allowbreak{}p5\_\allowbreak{}separated\_\allowbreak{}balanced & 2000 & 1000 & 0 & 52 \\
nt\_\allowbreak{}p5\_\allowbreak{}strong\_\allowbreak{}balanced & 500 & 1000 & 0 & 409 \\
nt\_\allowbreak{}p5\_\allowbreak{}strong\_\allowbreak{}balanced & 2000 & 1000 & 0 & 165 \\
ntsn\_\allowbreak{}p2\_\allowbreak{}moderate\_\allowbreak{}balanced & 1000 & 1000 & 8 & 147 \\
ntsn\_\allowbreak{}p5\_\allowbreak{}moderate\_\allowbreak{}balanced & 1000 & 1000 & 44 & 487 \\
\bottomrule
\end{longtable}
\endgroup

\begingroup\scriptsize\setlength{\tabcolsep}{2pt}
\begin{longtable}{p{0.22\linewidth}lrrrrrr}
\caption{Complete prespecified real-data screen. Dashes denote fits that failed before a valid covariance summary was obtained.}\label{tab:s-screen}\\
\toprule
Data & Family & $p$ & $n$ & ARI & Support & Inflation & Shrinkage\\
\midrule\endfirsthead
\toprule
Data & Family & $p$ & $n$ & ARI & Support & Inflation & Shrinkage\\
\midrule\endhead
breast\_\allowbreak{}cancer & Normal--Student-$t$ & 2 & 569 & 0.654 & 1 & 1.004 & 1.000 \\
breast\_\allowbreak{}cancer & Normal--skew-normal & 2 & 569 & 0.655 & 5 & 1.043 & 1.000 \\
breast\_\allowbreak{}cancer & Normal--Student-$t$ & 5 & 569 & 0.457 & 6 & 1.106 & 0.884 \\
breast\_\allowbreak{}cancer & Normal--skew-normal & 5 & 569 & 0.462 & 8 & 1.063 & 0.811 \\
wine & Normal--Student-$t$ & 2 & 178 & 0.389 & 2 & 1.418 & 0.799 \\
wine & Normal--skew-normal & 2 & 178 & 0.374 & 2 & 1.095 & 0.852 \\
wine & Normal--Student-$t$ & 5 & 178 & 0.359 & 1 & 1.024 & 1.000 \\
wine & Normal--skew-normal & 5 & 178 & 0.357 & 1 & 1.020 & 0.402 \\
digits\_\allowbreak{}0\_\allowbreak{}1 & Normal--Student-$t$ & 2 & 360 & 0.989 & 0 & 1.000 & 1.000 \\
digits\_\allowbreak{}0\_\allowbreak{}1 & Normal--skew-normal & 2 & 360 & 0.989 & 0 & 1.000 & 1.000 \\
digits\_\allowbreak{}0\_\allowbreak{}1 & Normal--Student-$t$ & 5 & 360 & 0.989 & 0 & 1.000 & 1.000 \\
digits\_\allowbreak{}0\_\allowbreak{}1 & Normal--skew-normal & 5 & 360 & -- & -- & -- & -- \\
dry\_\allowbreak{}bean\_\allowbreak{}cali\_\allowbreak{}barbunya & Normal--Student-$t$ & 2 & 2952 & 0.282 & 44 & 1.650 & 1.000 \\
dry\_\allowbreak{}bean\_\allowbreak{}cali\_\allowbreak{}barbunya & Normal--skew-normal & 2 & 2952 & 0.433 & 13 & 1.082 & 1.000 \\
dry\_\allowbreak{}bean\_\allowbreak{}cali\_\allowbreak{}barbunya & Normal--Student-$t$ & 5 & 2952 & 0.814 & 25 & 1.158 & 1.000 \\
dry\_\allowbreak{}bean\_\allowbreak{}cali\_\allowbreak{}barbunya & Normal--skew-normal & 5 & 2952 & 0.790 & 14 & 1.045 & 1.000 \\
dry\_\allowbreak{}bean\_\allowbreak{}dermason\_\allowbreak{}seker & Normal--Student-$t$ & 2 & 5573 & 0.867 & 33 & 1.284 & 1.000 \\
dry\_\allowbreak{}bean\_\allowbreak{}dermason\_\allowbreak{}seker & Normal--skew-normal & 2 & 5573 & 0.886 & 24 & 1.782 & 0.975 \\
dry\_\allowbreak{}bean\_\allowbreak{}dermason\_\allowbreak{}seker & Normal--Student-$t$ & 5 & 5573 & 0.825 & 21 & 1.180 & 0.664 \\
dry\_\allowbreak{}bean\_\allowbreak{}dermason\_\allowbreak{}seker & Normal--skew-normal & 5 & 5573 & -- & -- & -- & -- \\
dry\_\allowbreak{}bean\_\allowbreak{}horoz\_\allowbreak{}sira & Normal--Student-$t$ & 2 & 4564 & 0.908 & 12 & 1.040 & 1.000 \\
dry\_\allowbreak{}bean\_\allowbreak{}horoz\_\allowbreak{}sira & Normal--skew-normal & 2 & 4564 & 0.902 & 14 & 1.077 & 1.000 \\
dry\_\allowbreak{}bean\_\allowbreak{}horoz\_\allowbreak{}sira & Normal--Student-$t$ & 5 & 4564 & 0.863 & 12 & 1.071 & 1.000 \\
dry\_\allowbreak{}bean\_\allowbreak{}horoz\_\allowbreak{}sira & Normal--skew-normal & 5 & 4564 & -- & -- & -- & -- \\
\bottomrule
\end{longtable}
\endgroup

\begingroup\scriptsize\setlength{\tabcolsep}{2pt}
\begin{longtable}{p{0.30\linewidth}rrrrrr}
\caption{Parameter-level results for Dermason--Seker.}\label{tab:s-dermason-parameters}\\
\toprule
Coordinate & Estimate & Naive SE & Corrected SE & Local SD & Naive/SD & Corrected/SD\\
\midrule\endfirsthead
\toprule
Coordinate & Estimate & Naive SE & Corrected SE & Local SD & Naive/SD & Corrected/SD\\
\midrule\endhead
weight\_\allowbreak{}logit\_\allowbreak{}1\_\allowbreak{}vs\_\allowbreak{}2 & -0.642 & 0.028 & 0.037 & 0.036 & 0.789 & 1.030 \\
component\_\allowbreak{}1\_\allowbreak{}normal\_\allowbreak{}coordinate\_\allowbreak{}0 & 3.363 & 0.026 & 0.034 & 0.032 & 0.833 & 1.074 \\
component\_\allowbreak{}1\_\allowbreak{}normal\_\allowbreak{}coordinate\_\allowbreak{}1 & -0.621 & 0.046 & 0.060 & 0.061 & 0.748 & 0.988 \\
component\_\allowbreak{}1\_\allowbreak{}normal\_\allowbreak{}coordinate\_\allowbreak{}2 & 0.149 & 0.018 & 0.021 & 0.023 & 0.790 & 0.916 \\
component\_\allowbreak{}1\_\allowbreak{}normal\_\allowbreak{}coordinate\_\allowbreak{}3 & 0.599 & 0.047 & 0.060 & 0.058 & 0.807 & 1.032 \\
component\_\allowbreak{}1\_\allowbreak{}normal\_\allowbreak{}coordinate\_\allowbreak{}4 & 0.649 & 0.017 & 0.027 & 0.029 & 0.576 & 0.949 \\
component\_\allowbreak{}2\_\allowbreak{}student\_\allowbreak{}t\_\allowbreak{}coordinate\_\allowbreak{}0 & -1.755 & 0.028 & 0.039 & 0.036 & 0.783 & 1.075 \\
component\_\allowbreak{}2\_\allowbreak{}student\_\allowbreak{}t\_\allowbreak{}coordinate\_\allowbreak{}1 & 0.310 & 0.035 & 0.038 & 0.038 & 0.929 & 0.995 \\
component\_\allowbreak{}2\_\allowbreak{}student\_\allowbreak{}t\_\allowbreak{}coordinate\_\allowbreak{}2 & 0.483 & 0.013 & 0.019 & 0.017 & 0.736 & 1.082 \\
component\_\allowbreak{}2\_\allowbreak{}student\_\allowbreak{}t\_\allowbreak{}coordinate\_\allowbreak{}3 & 0.766 & 0.033 & 0.039 & 0.040 & 0.832 & 0.972 \\
component\_\allowbreak{}2\_\allowbreak{}student\_\allowbreak{}t\_\allowbreak{}coordinate\_\allowbreak{}4 & 0.632 & 0.014 & 0.015 & 0.015 & 0.939 & 0.986 \\
component\_\allowbreak{}2\_\allowbreak{}student\_\allowbreak{}t\_\allowbreak{}coordinate\_\allowbreak{}5 & 2.927 & 0.206 & 0.224 & 0.253 & 0.815 & 0.886 \\
\bottomrule
\end{longtable}
\endgroup

\begingroup\scriptsize\setlength{\tabcolsep}{2pt}
\begin{longtable}{p{0.30\linewidth}rrrrrr}
\caption{Parameter-level results for Cali--Barbunya.}\label{tab:s-cali-parameters}\\
\toprule
Coordinate & Estimate & Naive SE & Corrected SE & Local SD & Naive/SD & Corrected/SD\\
\midrule\endfirsthead
\toprule
Coordinate & Estimate & Naive SE & Corrected SE & Local SD & Naive/SD & Corrected/SD\\
\midrule\endhead
weight\_\allowbreak{}logit\_\allowbreak{}1\_\allowbreak{}vs\_\allowbreak{}2 & 1.187 & 0.043 & 0.113 & 0.150 & 0.290 & 0.753 \\
component\_\allowbreak{}1\_\allowbreak{}normal\_\allowbreak{}coordinate\_\allowbreak{}0 & 0.854 & 0.045 & 0.084 & 0.095 & 0.477 & 0.885 \\
component\_\allowbreak{}1\_\allowbreak{}normal\_\allowbreak{}coordinate\_\allowbreak{}1 & 0.673 & 0.035 & 0.055 & 0.068 & 0.510 & 0.817 \\
component\_\allowbreak{}1\_\allowbreak{}normal\_\allowbreak{}coordinate\_\allowbreak{}2 & 0.764 & 0.015 & 0.025 & 0.026 & 0.568 & 0.965 \\
component\_\allowbreak{}1\_\allowbreak{}normal\_\allowbreak{}coordinate\_\allowbreak{}3 & -0.720 & 0.034 & 0.051 & 0.060 & 0.565 & 0.846 \\
component\_\allowbreak{}1\_\allowbreak{}normal\_\allowbreak{}coordinate\_\allowbreak{}4 & 0.393 & 0.015 & 0.031 & 0.042 & 0.353 & 0.749 \\
component\_\allowbreak{}2\_\allowbreak{}student\_\allowbreak{}t\_\allowbreak{}coordinate\_\allowbreak{}0 & -2.687 & 0.097 & 0.171 & 0.221 & 0.438 & 0.776 \\
component\_\allowbreak{}2\_\allowbreak{}student\_\allowbreak{}t\_\allowbreak{}coordinate\_\allowbreak{}1 & -2.144 & 0.090 & 0.177 & 0.198 & 0.452 & 0.890 \\
component\_\allowbreak{}2\_\allowbreak{}student\_\allowbreak{}t\_\allowbreak{}coordinate\_\allowbreak{}2 & 0.776 & 0.035 & 0.043 & 0.050 & 0.697 & 0.855 \\
component\_\allowbreak{}2\_\allowbreak{}student\_\allowbreak{}t\_\allowbreak{}coordinate\_\allowbreak{}3 & -1.191 & 0.084 & 0.107 & 0.134 & 0.629 & 0.798 \\
component\_\allowbreak{}2\_\allowbreak{}student\_\allowbreak{}t\_\allowbreak{}coordinate\_\allowbreak{}4 & 0.525 & 0.041 & 0.051 & 0.056 & 0.723 & 0.898 \\
component\_\allowbreak{}2\_\allowbreak{}student\_\allowbreak{}t\_\allowbreak{}coordinate\_\allowbreak{}5 & 2.031 & 0.300 & 0.330 & 0.310 & 0.968 & 1.064 \\
\bottomrule
\end{longtable}
\endgroup


\section{Computational verification and reproducibility}

The implementation has unit tests for parameter packing and unpacking,
component log densities, family-assignment permutations, active-pair logic,
finite-difference gradients, curvature stabilization, and bootstrap failure
retention.  Vectorized contrast derivatives were compared with an independent
observation-wise implementation.  The skew-normal density used by fitting and
inference is identical, including its symmetric inverse square root.  Every
simulation seed is a deterministic function of the base seed, scenario,
sample size, and replication number.  Raw cluster-array chunks, standard
output and error logs, processed tables, and the scripts that generate this
supplement are retained in the reproduction project.
